% vim: tw=78 ai spell spelllang=cs
\documentclass{article}

%\usepackage{a4wide}
\usepackage{graphicx}
\usepackage[utf8x]{inputenc}
\usepackage[czech]{babel}
\usepackage[IL2]{fontenc}

\pagestyle{empty}

\begin{document}
\begin{flushright}
\today
\end{flushright}
\begin{center}
\Large\textbf{Posudek oponenta na diplomovou práci \\
\normalsize%
\textit{Bc. Matej Šuta:} String abstract domains}
\end{center}

Dle zadání měla diplomová práce doimplementovat do nástroje \emph{Canal}
další doménu pro abstraktní interpretaci, která bude schopna pracovat s
řetězci jazyka C ukončenými znakem \emph{NUL}. Text práce je na 33 stranách
strukturovaný do 6 kapitol. V úvodu autor krátce popisuje motivaci pro práci.
Následující kapitola popisuje různé metody analýzy kódu. Následuje kapitola s
teoretickými základy o abstraktní interpretaci a řetězcích, což je základ pro
další 2 kapitoly: návrh a implementaci řetězcových domén.  Práci ukončuje
závěr obsahující i krátký popis práce do budoucna.  \textbf{Rozsah} uvedených
kapitol je z pohledu diplomové práce i zadání obvyklý.

\textbf{Obsahově} dle mého názoru dostatečně popisuje teoretické základy
práce, návrh i implementaci. Nicméně z práce není úplně zřejmá její motivace.
Mělo být více rozvedeno, jaké má práce využití. Chybí také výsledky na nějakém
reálném softwaru. Jednak k ověření, že implementace opravdu funguje podle
návrhu, jednak k ilustraci její výkonnosti. Dále se v práci objevuje
nesrozumitelný (,,Most of the languages'' v kap. 2.2 odst. 2. -- kterých
jazyků?, ,,deal with them once they do'', str. 6 odst. 2: ,,implemented by a
programmer''? atd.) a nadbytečný text (zmínka o \textsc{Divine} v 2.3.1, s tou
informací se nijak nepracuje). Také je možné narazit na pasáže s nepřesnostmi,
např.:
\begin{itemize}
\item V 2.2: \texttt{grep} vrací spoustu falešných pozitiv. To nemusí být při
dobře zvoleném regulárním výrazu pravda.
\item V 2.3.1, odst. 1: často se stává, že nedostaneme výsledek žádný -- právě
kvůli explozi.
\item V 2.3.4: proč musí měnit \emph{každá} instrukce stav?
\item Definice 4.3 a 4.7: nesprávně použitý symbol $\leq$ namísto
$\sqsubseteq$.
\item Str. 30 nahoře: \textsc{Klee} není jediný symbolický exekutor nad
\textsc{Llvm} (nesprávně užito ,,namely''?).
\end{itemize}

Po \textbf{jazykové} stránce je práce na průměrné úrovni. Je zde několik
překlepů (,,is has'', ,,isBottom'' namísto ,,isTop'' na str. 32, dva otazníky
v 4.4), amerikanismů (,,behavior''), nesprávně použití množného čísla
(,,information'' je v tomto kontextu nepočitatelné) a spoustu chybějících
členů (např. ,,how program manipulates''). Občas dochází k chybnému dělení
slov (,,anal-ysis'' na str. 29). Před odkazy na poznámku pod čarou jsou
nesprávně mezery, nadpisy jsou malými písmeny. Kotvy odkazů u tabulek je
vhodnější vkládat na začátek tabulky, aby prohlížeč zobrazil tabulku celou.
Oceňuji volbu anglického jazyka.

\textbf{Praktická část} práce je v přiloženém \texttt{zip} souboru. Jelikož se
jedná o rozšíření již existujícího nástroje \emph{Canal}, je v práci přesně
uvedeno, co přesně bylo změněno a přidáno. Tento nově přidaný kód je čitelný a
jeví se být funkční.

Na závěr konstatuji, že práce \textbf{splňuje zadání a požadavky na
diplomovou práci}. Navrhuji práci \textbf{uznat} jako diplomovou a vzhledem
k výše uvedeným okolnostem doporučuji hodnocení stupněm B, uvede-li autor v
prezentaci obsáhlejší motivaci (ideálně případ užití) a výsledky měření
výkonnosti (např. na \textsc{Coreutils}), jinak C.

\vspace{3em}
\begin{flushright}
Jiří Slabý
\end{flushright}

\end{document}
